\customchapter{INTRODUCCIÓN}

El contenido audiovisual crece más rápido de lo que una persona puede ver, y para indexarlo,
previsualizarlo o navegarlo hay que saber dónde termina cada unidad de sentido y qué
segmentos, aunque estén separados en el tiempo, tratan el mismo tema. La \emph{segmentación
temática} agrupa el contenido en tópicos coherentes. A diferencia de la segmentación en
escenas, admite que un tema se abandone y se retome más adelante, y permite que el mismo
video se consuma de formas distintas según el interés y el tiempo de quien lo ve.

Esta tesis continúa dos trabajos previos del mismo grupo. El primero, A2Summ-X, extendió el
modelo de resumen multimodal A2Summ \cite{he2023a2summ} y se evaluó con 5 splits y 5 semillas
por configuración y pruebas t pareadas sobre SumMe y TVSum \cite{gygli2014summe,
song2015tvsum}. La rama de texto aportó una mejora pequeña ($\Delta$F1 $= +0.033$,
$p = 0.014$ en TVSum), los backbones visuales modernos no mejoraron el resultado y en
conjuntos de pocos videos la varianza entre corridas dominó sobre los efectos. El segundo, SegLoop,
usó un modelo visión--lenguaje (VLM) congelado \cite{bai2025qwen25vl} como juez para generar
pseudoetiquetas con las que se entrena un segmentador de escenas ligero \cite{mun2022bassl,
huang2020movienet}. De SegLoop queda una infraestructura verificada de extracción de señales
(reconocimiento de voz, descriptores visuales y detección de cortes) que CAVA reutiliza.

El análisis que cerró A2Summ-X encontró el problema que esta tesis aborda: al fusionar las
señales en un embedding o un score único, unas reescriben a otras antes de la decisión. En
A2Summ-X, la alineación inyecta la transcripción en la consulta visual, la compuerta softmax
reparte un peso total fijo entre modalidades y la cabeza de decisión solo ve la mezcla.
Cuatro mediciones independientes coincidieron: pesos de la compuerta de texto entre 0.75 y
0.89, colapso de la señal visual en BLiSS, correlación de importancia $\rho \approx 0$ fuera
de dominio y secuencias sin cortes detectables. De ahí sale la regla de diseño de CAVA:
mantener separada la identidad de cada señal hasta el momento de decidir.

En CAVA, cada señal (transcripción, imagen, voz y las descripciones que un VLM produce por
bloques del video) calcula sus propias fronteras con su propio percentil y asigna sus propias
etiquetas de tópico. El tópico final surge de un consenso de votos con quórum. Una señal sin
evidencia en un tramo se abstiene, y cada grupo resultante registra qué señales lo sostienen
(su \emph{procedencia}). Sobre este gestor se construye el mapa de tópicos del video, con
enlaces entre segmentos no adyacentes, dos niveles de granularidad y contexto por tramo. Del
mapa se derivan cinco formatos de versión: capítulos, recapitulación de 60 segundos,
destacados del 10\,\%, hilos de tópico e hilos por hablante. El documento reporta el estado
del sistema al congelarlo: el gestor está validado en el split de test de MovieNet, las
compuertas de coherencia y de versiones están medidas, tres compuertas no se cumplen y la
evaluación con personas (EXP-040) está diseñada.

\introsection{Formulación del problema}

¿Es posible segmentar temáticamente video largo de contenido heterogéneo (película,
documental, reunión, video web) y producir versiones adaptadas al consumidor, sin anotación
humana de tópicos y sin fusionar representaciones entre señales, de modo que:
(i) el mapa de tópicos por consenso supere a la fusión de señales y a la segmentación
contigua en una métrica pre-registrada sobre datos no vistos;
(ii) los tópicos sean coherentes dentro de cada señal (margen intra$-$inter positivo);
(iii) cada versión respete su presupuesto de duración; y
(iv) un juez independiente reconozca los enlaces no contiguos como retornos del mismo tema?

\introsection{Objetivos de investigación}

\paragraph{Objetivo general}
Construir y validar experimentalmente CAVA, un sistema que segmenta temáticamente video largo
mediante agregación multimodal por consenso, sin fusión de representaciones, con abstención y
procedencia por señal, y que produce varias versiones del contenido adaptadas al consumidor,
validadas con compuertas pre-registradas sobre datos no vistos.

\paragraph{Objetivos específicos}
\begin{enumerate}
  \item Implementar las señales del sistema sobre la infraestructura de SegLoop:
        transcripción \cite{radford2023whisper} proyectada a embeddings, señal visual
        SigLIP \cite{zhai2023siglip} por unidad, voz (log-mel y tono) y descripciones de un
        VLM congelado \cite{bai2025qwen25vl}, con perfiles por tipo de contenido (película,
        documental, reunión) y ajuste de fronteras a cortes visuales.
  \item Implementar el gestor multimodal sin fusión: fronteras y etiquetas propias por
        señal, consenso de cortes por votos ponderados con tolerancia, enlace de tópicos por
        quórum (componentes conexas o enlace promedio), abstención de señales sin evidencia y
        procedencia por señal en cada grupo.
  \item Producir el mapa de tópicos con retornos, dos niveles de granularidad, etiquetas
        generadas por un modelo de lenguaje, nombres de personajes propagados por los
        enlaces y contexto por tramo, junto con la identidad de hablante (embeddings
        WavLM-sv \cite{chen2022wavlm} agrupados con Ward y BIC).
  \item Generar los cinco formatos de versión con verificación automática de presupuestos
        (compuerta G-versión).
  \item Validar el sistema con compuertas pre-registradas: el gestor frente a la fusión y al
        gestor contiguo en el split de test de MovieNet (47 películas); la coherencia
        intra$-$inter por señal; la precisión de los enlaces no contiguos con un juez ciego;
        la identidad de hablante en AMI \cite{carletta2007ami} y VoxConverse
        \cite{chung2020voxconverse}; y la calidad de los clips frente a referencias
        automáticas y bases simples.
\end{enumerate}

\introsection{Justificación}

La segmentación temática con enlaces no contiguos es necesaria para aplicaciones que la
segmentación en escenas no cubre: navegar un video por tema, generar recapitulaciones con un
tiempo fijo y reunir en un hilo todas las apariciones de un tema o de un hablante. Los
sistemas actuales dependen de VLM costosos en inferencia \cite{berman2025scenevlm} o
fusionan las modalidades antes de decidir, y en ese segundo caso el trabajo previo del grupo
mostró que una señal termina dominando a las demás. Si el consenso tardío produce mapas de
tópicos iguales o mejores que la fusión, se gana además la posibilidad de saber qué señales
sostienen cada grupo, algo que una representación fusionada no permite recuperar.

En lo metodológico, el trabajo aplica a un problema nuevo el esquema de compuertas
pre-registradas del grupo: separación de desarrollo, confirmación y test, y selección de
hiperparámetros sin consultar test. Como producto queda un pipeline reproducible (115
corridas con manifiesto que registra que no se usó anotación, 25 pruebas automatizadas y un
registro de decisiones de arquitectura) y un inventario del estado de validación de cada
módulo por tipo de contenido.

\introsection{Alcance y limitaciones}

La investigación cubre el ciclo completo, desde la extracción de señales hasta las versiones
por consumidor, con estos límites:
(i) los dominios evaluados son película (MovieNet), documental (BBC Planet Earth), reunión
(AMI) y video web y de usuario (VidChapters \cite{ventura2024vidchapters}, TVSum y SumMe),
sin anotación nueva de tópicos salvo la del conjunto de la evaluación con personas;
(ii) la calidad percibida de las versiones se mide con presupuestos y referencias
automáticas, y la preferencia humana (EXP-040) queda diseñada pero no ejecutada;
(iii) las identidades por apariencia (P\#) no se publican como identidad, y solo la identidad
por voz está validada, y únicamente en habla conversacional;
(iv) el texto en pantalla (OCR) y el espectro de audio están definidos como contratos de
señal pero no se extraen; y
(v) el cómputo se limita al clúster institucional (GPU RTX A6000 de 48\,GB y A100),
con límites de tiempo por trabajo.
